\documentclass[11pt]{article}
% Shared preamble for the rigor documents (Lamport-style structured proofs).  \input this file after \documentclass.
\usepackage[T1]{fontenc}\usepackage{lmodern}\usepackage{microtype}
\usepackage[margin=1in]{geometry}
\usepackage{amsmath,amssymb,amsthm,mathtools}
\usepackage{xcolor}\usepackage{enumitem}\usepackage{booktabs}\usepackage{graphicx}\usepackage{hyperref}
\usepackage[most]{tcolorbox}
\definecolor{navy}{HTML}{17365D}\definecolor{teal}{HTML}{117A8B}\definecolor{warn}{HTML}{A33A2B}\definecolor{okgreen}{HTML}{2E7D4F}
\hypersetup{colorlinks=true,linkcolor=navy,citecolor=teal,urlcolor=teal}
\theoremstyle{plain}
\newtheorem{theorem}{Theorem}[section]\newtheorem{lemma}[theorem]{Lemma}\newtheorem{proposition}[theorem]{Proposition}\newtheorem{corollary}[theorem]{Corollary}
\theoremstyle{definition}\newtheorem{definition}[theorem]{Definition}\newtheorem{remark}[theorem]{Remark}\newtheorem{claim}[theorem]{Claim}\newtheorem{issue}[theorem]{Issue}
% ---------- Lamport structured proofs ----------
% \lstep{level}{number}{assertion}   -> indented "<level>number. assertion"
% \lproof                             -> "PROOF:" line (start of the sub-proof of the preceding step)
% \lby{justification}                 -> "BY ..." leaf justification for the preceding step
% \lqed{level}{number}                -> QED step of a sub-proof
% keywords: \lassume \lprove \lsuffices \lcase \ldefine \lpick \lnew
\newlength{\lindent}\setlength{\lindent}{1.7em}
\newcommand{\lstep}[3]{\par\smallskip\noindent\hspace*{\dimexpr\numexpr#1-1\relax\lindent\relax}{\color{navy}$\langle#1\rangle#2$.}\ #3}
\newcommand{\lproof}{\par\noindent\hspace*{\lindent}{\scshape Proof:}}
\newcommand{\lby}[1]{\par\noindent\hspace*{\lindent}{\scshape By}\ #1.}
\newcommand{\lqed}[2]{\par\smallskip\noindent\hspace*{\dimexpr\numexpr#1-1\relax\lindent\relax}{\color{navy}$\langle#1\rangle#2$.}\ {\scshape Q.E.D.}}
\newcommand{\lassume}{{\scshape Assume}\ }\newcommand{\lprove}{{\scshape Prove}\ }\newcommand{\lsuffices}{{\scshape Suffices}\ }
\newcommand{\lcase}{{\scshape Case}\ }\newcommand{\ldefine}{{\scshape Define}\ }\newcommand{\lpick}{{\scshape Pick}\ }\newcommand{\lnew}{{\scshape New}\ }
% ---------- verdict boxes ----------
\newtcolorbox{verdictok}{colback=okgreen!8,colframe=okgreen,title={\bfseries Verdict: verified},fonttitle=\small}
\newtcolorbox{verdictgap}{colback=orange!10,colframe=orange!80!black,title={\bfseries Verdict: gap / caveat},fonttitle=\small}
\newtcolorbox{verdictbreak}{colback=warn!10,colframe=warn,title={\bfseries Verdict: proof-breaking issue},fonttitle=\small}
\newtcolorbox{citedbox}[1]{colback=navy!5,colframe=navy!60,title={\bfseries Cited result: #1},fonttitle=\small,breakable}
\setlength{\parindent}{0pt}\setlength{\parskip}{4pt}\allowdisplaybreaks
\newcommand{\cA}{\mathcal A}\newcommand{\cF}{\mathcal F}\newcommand{\cM}{\mathfrak M}\newcommand{\cN}{\mathfrak N}

% extra calligraphic shortcuts (\providecommand: no clash if a document defines its own)
\providecommand{\cB}{\mathcal B}\providecommand{\cC}{\mathcal C}\providecommand{\cD}{\mathcal D}\providecommand{\cE}{\mathcal E}\providecommand{\cG}{\mathcal G}\providecommand{\cH}{\mathcal H}\providecommand{\cI}{\mathcal I}\providecommand{\cJ}{\mathcal J}\providecommand{\cK}{\mathcal K}\providecommand{\cL}{\mathcal L}\providecommand{\cO}{\mathcal O}\providecommand{\cP}{\mathcal P}\providecommand{\cQ}{\mathcal Q}\providecommand{\cR}{\mathcal R}\providecommand{\cS}{\mathcal S}\providecommand{\cT}{\mathcal T}\providecommand{\cU}{\mathcal U}\providecommand{\cV}{\mathcal V}\providecommand{\cW}{\mathcal W}\providecommand{\cX}{\mathcal X}\providecommand{\cY}{\mathcal Y}\providecommand{\cZ}{\mathcal Z}
\newcommand{\Fid}{\operatorname{F}}\newcommand{\Tr}{\operatorname{Tr}}\newcommand{\eps}{\varepsilon}\newcommand{\dd}{\,\mathrm d}

\title{Exterior extensions and the optimised Uhlmann bound:\\ numerics in the spectral NS-circle model (track E2)}
\author{E2}
\date{2026-09-15}
\begin{document}
\maketitle
\tableofcontents

\section{Setup, notation, conventions}\label{sec:setup}

All numerics are in \texttt{cft\_cmi/numerics/optimality\_all/}: the module
\texttt{e2\_common.py} and the scripts \texttt{e2\_subspaces.py},
\texttt{e2\_dist\_ub.py}, \texttt{e2\_exact.py}, with outputs \texttt{e2\_*.out}.

\subsection{The model}
$L$ sites on the circle, $\theta_j=2\pi j/L$ folded to $(-\pi,\pi]$, NS (antiperiodic)
momenta $\nu=n+\tfrac12$; $P$ is the exact Hardy projection $\mathbf 1_{\nu>0}$ conjugated
by the DFT, so $P=P^*=P^2$ and $k:=\operatorname{rank}P=L/2$.  The arcs are
$D=(0,\theta_D)$, $\theta_D=2\arctan L_g$, $A=(\theta_A,0)$, $\theta_A=-2\arctan a$,
$I=A\cup D$, and we use $a=2$, $L_g=1$ throughout, as in \texttt{s10\_theta.py};
$m:=|I|$, $m_c:=|I^c|=L-m$.  The compression generator is
$D_\chi=\tfrac12(\chi\,\partial+\partial\,\chi)$ (anti-self-adjoint) with the UV-tapered
$\partial$ of \texttt{s10\_theta.py} (taper window $(0.30,0.42)$) and
$\chi=\cos\theta-1$ on $D$, continued by $(\cos t-1)e^{-3(t-\theta_D)^2}$ on $I^c$;
$M:=P^\perp D_\chi P$.  $G_*$ is the S10 rotation direction obtained by conjugate
gradients, $v:=P^\perp G_*P$, and $\theta$ the relative gain.

\subsection{The real Hilbert space and the three subspaces}
$V:=\{A:\ A=P^\perp AP\}$ with the \emph{real} inner product
$\langle A,B\rangle:=\Re\Tr(A^*B)$; $\dim_{\mathbb R}V=2k(L-k)=L^2/2$.  With
$E_I,E_c:=E_{I^c}$ the coordinate projections,
\begin{align}
 \mathcal N&:=\operatorname{span}_{\mathbb R}\{P^\perp\ell P:\ \ell=\ell^*=E_I\ell E_I\},&
 \mathrm{Ext}&:=\operatorname{span}_{\mathbb R}\{P^\perp Z P:\ Z^*=-Z=E_cZE_c\}.
\end{align}
$\Pi$ is the real-orthogonal projection onto $\mathcal N$ and $u=\Pi M$.
For $A\in V$ one has the two pairings
\begin{equation}\label{eq:pair}
 \langle P^\perp\ell P,A\rangle=\tfrac12\Tr\bigl(\ell\,E_I(A+A^*)E_I\bigr),\qquad
 \langle P^\perp Z P,A\rangle=-\tfrac12\Tr\bigl(Z\,E_c(A-A^*)E_c\bigr),
\end{equation}
whence $\mathcal N^\perp=\{A\in V: E_I(A+A^*)E_I=0\}$ and
$\mathrm{Ext}^\perp=\{A\in V:E_c(A-A^*)E_c=0\}$; in particular
$\mathrm{Ext}=\mathcal N^\perp$ iff $\mathrm{Ext}^\perp=\mathcal N$.

\subsection{What is actually free: the column, not the block}\label{ssec:column}
The brief writes the admissible generators as $Z=Z_I\oplus Z_e$, block diagonal for
$L^2(I)\oplus L^2(I^c)$.  That parametrisation is correct only if $\widetilde W$
preserves $L^2(I)$, which the compression flow does \emph{not}: $\chi(\theta_D)=\cos\theta_D-1=-1\neq0$,
so the flow pushes points across $\partial I$ and $W_I$ is an isometry of $L^2(I)$ that is
not onto.  The correct linearisation is: prescribing $\widetilde W|_{L^2(I)}=W_I$ fixes the
whole \emph{column} $ZE_I$ (both $E_IZE_I$ and $E_cZE_I$) and leaves exactly the block
$E_cZE_c$ free.  Since $P^\perp(E_cZE_c)P\in\mathrm{Ext}$ by definition, and since
$E_cD_\chi E_c$ is itself an admissible $Z_e$, one may equivalently keep the full $D_\chi$
and optimise over all anti-Hermitian $Z_e$ supported on $I^c$:
\begin{equation}\label{eq:UB}
 \mathrm{UB}(X):=\min_{Z_e}\tfrac12\bigl\|X+P^\perp Z_eP\bigr\|_2^2
 =\tfrac12\operatorname{dist}(X,\mathrm{Ext})^2
 =\tfrac12\Bigl(\|\Pi X\|^2+\operatorname{dist}\bigl((1-\Pi)X,\mathrm{Ext}\bigr)^2\Bigr),
\end{equation}
the last equality because $\mathrm{Ext}\subseteq\mathcal N^\perp$ (L1), so
$\mathrm{Proj}_{\mathrm{Ext}}\Pi X=0$.  This is what the scripts compute.  The naive matrix
truncation $Z_I=E_ID_\chi E_I$ is \emph{not} admissible; see \S\ref{sec:dist}.

\subsection{Exact diagonalisation of both Gram forms}
Because $\ell=E_I\ell E_I$ and $Z=E_cZE_c$, the two compressions are exact:
with $Q:=E_IPE_I$ (eigenvalues $q_i$, eigenbasis $U$) and $Q_c:=E_cPE_c$
(eigenvalues $q'_i$, eigenbasis $U_c$),
\begin{equation}\label{eq:gram}
 E_I(P^\perp\ell P+P\ell P^\perp)E_I=(1-Q)\ell Q+Q\ell(1-Q),\qquad
 E_c(P^\perp ZP+PZP^\perp)E_c=(1-Q_c)ZQ_c+Q_cZ(1-Q_c).
\end{equation}
In the respective eigenbases these are multiplication by
$d_{ij}:=q_i(1-q_j)+q_j(1-q_i)$ and $d^c_{ij}:=q'_i(1-q'_j)+q'_j(1-q'_i)$, so
\begin{equation}\label{eq:closed}
 \|P^\perp\ell P\|_2^2=\tfrac12\sum_{ij}d_{ij}|\ell_{ij}|^2,\qquad
 \|\Pi A\|^2=\tfrac12\sum_{ij}\frac{|\eta_{ij}|^2}{d_{ij}},\qquad
 \|\mathrm{Proj}_{\mathrm{Ext}}A\|^2=\tfrac12\sum_{ij}\frac{|K_{ij}|^2}{d^c_{ij}},
\end{equation}
with $\eta:=E_I(A+A^*)E_I$ in the $Q$-basis and $K:=E_c(A-A^*)E_c$ in the $Q_c$-basis.
Both projections and the minimiser $Z_e=-K/(d^c+\mu)$ (Tikhonov parameter $\mu\ge0$
bounding $\|Z_e\|_2$) are therefore available in closed form at $O(L^3)$ cost --- no
large least-squares system is ever assembled.  The first formula in \eqref{eq:closed}
reproduces \texttt{s10\_theta.py}'s \texttt{riesz()} and its weight $w=1/d$.


\section{Dimensions of $\mathcal N$, $\mathcal N^\perp$, $\mathrm{Ext}$}\label{sec:dims}

\begin{lemma}[Exact dimension count in the finite model]\label{lem:dim}
Let $m=|I|>k=L/2$ and $z:=m-k$.  Generically (i.e.\ when the $m\times k$ matrices
$E_IV_+$ and $E_IV_-$, $V_\pm$ the isometries onto $\operatorname{ran}P^{(\perp)}$, have
full rank $k$, and $E_cV_\pm$ have full rank $m_c$), 
\begin{equation}\label{eq:dimid}
 \dim_{\mathbb R}\mathcal N=m^2-2z^2,\qquad
 \dim_{\mathbb R}\mathcal N^\perp=\tfrac{L^2}2-m^2+2z^2=(k-z)^2=(L-m)^2
 =\dim_{\mathbb R}\mathrm{Ext},
\end{equation}
and consequently $\mathrm{Ext}=\mathcal N^\perp$ exactly.
\end{lemma}
\begin{proof}
\lstep{1}{1}{$Q=E_IPE_I$ has exactly $z$ eigenvalues $0$ and exactly $z$ eigenvalues $1$,
 and no others in $\{0,1\}$.}
\lproof
\lstep{2}{1}{$Q=(E_IV_+)(E_IV_+)^*$ with $E_IV_+$ of size $m\times k$ and rank $k$, so
 $\dim\ker Q=m-k=z$; the same for $1-Q=(E_IV_-)(E_IV_-)^*$.}
\lqed{2}{2}\lby{$\langle2\rangle1$}
\lstep{1}{2}{$\dim_{\mathbb R}\mathcal N=\#\{(i,j):d_{ij}>0\}=m^2-2z^2$.}
\lby{the first identity of \eqref{eq:closed}: the real-linear map $\ell\mapsto P^\perp\ell P$
 is, in the $Q$-eigenbasis, multiplication of the $m^2$ real coordinates of $\ell$ by
 $\sqrt{d_{ij}/2}$, and $d_{ij}=0$ iff $q_i=q_j=0$ or $q_i=q_j=1$, which by
 $\langle1\rangle1$ happens for exactly $z^2+z^2$ ordered pairs}
\lstep{1}{3}{$\dim_{\mathbb R}\mathrm{Ext}=m_c^2$.}
\lby{the same argument with $Q_c$; here $m_c\le k$, so $E_cV_\pm$ have full rank $m_c$,
 no eigenvalue of $Q_c$ equals $0$ or $1$, and all $d^c_{ij}>0$}
\lstep{1}{4}{$\dim_{\mathbb R}V=2k(L-k)=L^2/2$, hence
 $\dim\mathcal N^\perp=\tfrac{L^2}2-m^2+2z^2$; with $m=k+z$ and $k=L/2$ this equals
 $k^2-2kz+z^2=(k-z)^2=(L-m)^2=m_c^2$.}
\lby{$\langle1\rangle2$ and algebra}
\lstep{1}{5}{$\mathrm{Ext}\subseteq\mathcal N^\perp$ and the two have the same finite
 dimension, hence they are equal.}
\lby{(L1) of the Phase 4 brief, restated as \eqref{eq:pair}, $\langle1\rangle3$, $\langle1\rangle4$}
\lqed{1}{6}
\end{proof}

\subsection{Numerical verification}
Table~\ref{tab:dim} gives the exact prediction and the two independent numerical checks.
\emph{(A) Brute force} (\texttt{e2\_subspaces.py}, part A): explicit real bases of
$\mathcal N$ and $\mathrm{Ext}$ in the coordinates $A\mapsto V_-^*AV_+\in\mathbb C^{k\times k}$,
ranks by SVD at relative tolerance $10^{-10}$, and the principal angles between
$\mathrm{Ext}$ and $\mathcal N^\perp$.  \emph{(B) Closed form} (part B): the counts
$\#\{d_{ij}>\tau\,d_{\max}\}$ from \eqref{eq:closed}.

\begin{table}[htbp]\centering\small
\begin{tabular}{rrrr rr rr l}
\toprule
$L$&$m$&$m_c$&$\dim V$&$\dim\mathcal N$&$\dim\mathcal N^\perp$&rank$_{\rm num}\mathcal N$&rank$_{\rm num}$Ext&$\cos\angle(\mathrm{Ext},\mathcal N^\perp)$\\
\midrule
32&18&14& 512& 316& 196& 316& 196& $1.0000000000$ (all)\\
48&27&21&1152& 711& 441& 711& 441& $1.0000000000$ (all)\\
64&37&27&2048&1319& 729&1297& 727& $1.0000000000$ (all)\\
\midrule
96 & 56& 40& 4608& 3008& 1600&3028&1568&---\\
128& 76& 52& 8192& 5488& 2704&5325&2559&---\\
192&114& 78&18432&12348& 6084&11892&5725&---\\
256&153&103&32768&22159&10609&21301&9809&---\\
384&230&154&73728&50012&23716&48851&22489&---\\
\bottomrule
\end{tabular}
\caption{Dimensions.  Columns 5--6 are the exact prediction \eqref{eq:dimid}; columns 7--8
the numerical ranks (rows $L\le64$: SVD of explicit bases at relative tolerance
$10^{-10}$; rows $L\ge96$: counts of $d_{ij},d^c_{ij}$ above $\tau=10^{-8}$ times the
maximum).  Last column: all principal angles between $\mathrm{Ext}$ and
$\mathcal N^\perp$ vanish.  The deficits at $L\ge64$ are pure roundoff: they grow with
$L$ and shrink when $\tau$ is lowered, and they never produce a nonzero principal angle.}
\label{tab:dim}
\end{table}

\begin{verdictok}
$\mathrm{Ext}=\mathcal N^\perp$ holds \emph{exactly} in the finite NS-circle model, by
Lemma~\ref{lem:dim} and by two independent numerical routes.  The answer to question (Q)
of the Phase 4 brief is YES at every finite $L$; the brief's dimension count
$\dim\mathcal N^\perp=2k(n-k)-m^2\le(n-m)^2$ ``with equality iff $m=n/2$'' is not quite
right --- the correct statement is \eqref{eq:dimid}, and equality holds for \emph{every}
$m$, the apparent defect $m^2-\dim\mathcal N=2z^2$ being exactly the rank deficiency of
$\ell\mapsto P^\perp\ell P$ forced by $\operatorname{rank}Q=\operatorname{rank}(1-Q)=k$.
\end{verdictok}

\subsection{Conditioning: what the continuum limit could still spoil}
$\mathrm{Ext}=\mathcal N^\perp$ as sets says nothing about the \emph{cost} of reaching a
given direction: the gains of $Z\mapsto P^\perp ZP$ are $\sqrt{d^c_{ij}/2}\le1/\sqrt2$ and
they are spread over many decades.  Quantiles of $\log_{10}$(gain) for $\mathrm{Ext}$
(min, 5\%, 25\%, 50\%, 75\%, max), from \texttt{e2\_subspaces.out}:
\begin{center}\small
\begin{tabular}{rcccccc}
\toprule
$L$&min&5\%&25\%&50\%&75\%&max\\\midrule
96 &$-8.13$&$-7.40$&$-2.54$&$-0.30$&$-0.15$&$-0.15$\\
128&$-8.13$&$-7.83$&$-2.93$&$-0.30$&$-0.15$&$-0.15$\\
192&$-9.04$&$-7.81$&$-4.31$&$-0.30$&$-0.15$&$-0.15$\\
256&$-8.35$&$-7.96$&$-5.41$&$-0.17$&$-0.15$&$-0.15$\\
384&$-9.11$&$-7.89$&$-7.51$&$-0.20$&$-0.15$&$-0.15$\\
\bottomrule
\end{tabular}
\end{center}
The lower quartile falls from $10^{-2.5}$ to $10^{-7.5}$ between $L=96$ and $L=384$: a
growing fraction of $\mathcal N^\perp$ costs an exponentially large $\|Z_e\|_2$.  This is
exactly the shape of ``dense but not closed'' in the continuum, and it is the reason the
set-theoretic answer of Lemma~\ref{lem:dim} does not by itself settle (UH).  What settles
it is that the \emph{specific} direction $(1-\Pi)M$ turns out to sit in the well-conditioned
part; see \S\ref{sec:dist} and \S\ref{sec:ub}.


\section{Invisible components and their distance to $\mathrm{Ext}$}\label{sec:dist}

Script \texttt{e2\_dist\_ub.py}, output \texttt{e2\_dist\_ub.out}.  Because
$\mathrm{Ext}\subseteq\mathcal N^\perp$, $\mathrm{Proj}_{\mathrm{Ext}}\Pi X=0$ and
\begin{equation}\label{eq:d2}
 \operatorname{dist}\bigl((1-\Pi)X,\mathrm{Ext}\bigr)^2
 =\|X\|^2-\|\Pi X\|^2-\|\mathrm{Proj}_{\mathrm{Ext}}X\|^2 ,
\end{equation}
all three terms given in closed form by \eqref{eq:closed}.  Every entry below was also
checked for consistency by the ``leak'' diagnostics: the weight of $\eta$ on the
$d_{ij}=0$ modes and of $K$ on the $d^c_{ij}=0$ modes, which must vanish identically
(otherwise the pairing \eqref{eq:pair} would be unbounded); these are $\le10^{-27}$
everywhere.

\begin{table}[htbp]\centering\small
\begin{tabular}{r l rrr r r}
\toprule
$L$&$X$&$\|X\|^2$&$\|\Pi X\|^2$&$\|(1-\Pi)X\|^2$&$\operatorname{dist}((1-\Pi)X,\mathrm{Ext})^2$&rel.\ dist.\\
\midrule
128&$M$          &$3.21604\mathrm{e}{-2}$&$7.35861\mathrm{e}{-3}$&$2.48018\mathrm{e}{-2}$&$+3.5\mathrm{e}{-12}$&$1.2\mathrm{e}{-5}$\\
128&$v$          &$3.40538\mathrm{e}{-3}$&$2.14244\mathrm{e}{-3}$&$1.26294\mathrm{e}{-3}$&$+6.1\mathrm{e}{-18}$&$6.9\mathrm{e}{-8}$\\
128&$M+v$        &$3.06961\mathrm{e}{-2}$&$5.21616\mathrm{e}{-3}$&$2.54799\mathrm{e}{-2}$&$+3.5\mathrm{e}{-12}$&$1.2\mathrm{e}{-5}$\\
192&$M$          &$3.19940\mathrm{e}{-2}$&$7.27173\mathrm{e}{-3}$&$2.47223\mathrm{e}{-2}$&$-4.3\mathrm{e}{-11}$&$0$\\
192&$v$          &$3.51824\mathrm{e}{-3}$&$2.09971\mathrm{e}{-3}$&$1.41853\mathrm{e}{-3}$&$-1.0\mathrm{e}{-16}$&$0$\\
192&$M+v$        &$3.04094\mathrm{e}{-2}$&$5.17202\mathrm{e}{-3}$&$2.52374\mathrm{e}{-2}$&$-4.3\mathrm{e}{-11}$&$0$\\
256&$M$          &$3.19400\mathrm{e}{-2}$&$7.38711\mathrm{e}{-3}$&$2.45529\mathrm{e}{-2}$&$-2.0\mathrm{e}{-12}$&$0$\\
256&$v$          &$3.69905\mathrm{e}{-3}$&$2.13796\mathrm{e}{-3}$&$1.56109\mathrm{e}{-3}$&$-8.3\mathrm{e}{-17}$&$0$\\
256&$M+v$        &$3.03221\mathrm{e}{-2}$&$5.24916\mathrm{e}{-3}$&$2.50729\mathrm{e}{-2}$&$-2.0\mathrm{e}{-12}$&$0$\\
384&$M$          &$3.19026\mathrm{e}{-2}$&$7.41022\mathrm{e}{-3}$&$2.44924\mathrm{e}{-2}$&$-1.3\mathrm{e}{-10}$&$0$\\
384&$v$          &$3.91398\mathrm{e}{-3}$&$2.16633\mathrm{e}{-3}$&$1.74765\mathrm{e}{-3}$&$-1.5\mathrm{e}{-16}$&$0$\\
384&$M+v$        &$3.02638\mathrm{e}{-2}$&$5.24389\mathrm{e}{-3}$&$2.50200\mathrm{e}{-2}$&$-1.3\mathrm{e}{-10}$&$0$\\
\bottomrule
\end{tabular}
\caption{Invisible components and their distance to $\mathrm{Ext}$.  $M=P^\perp D_\chi P$,
$v=P^\perp G_*P$ with $G_*=\eps_*\,x$ the CG direction of S10 ($\eps_*=-1.0000$ in units
of $s$, as predicted by $\langle u,v_x\rangle=\|v_x\|^2=\langle b,x\rangle$).  The
distances are at the roundoff level --- they are of the size of the accumulated error in a
difference of two numbers of size $10^{-2}$ computed through a $\sum|{\cdot}|^2/d$ with
$d$ down to $10^{-17}$, and they change sign with $L$.  ``rel.\ dist.''\ is
$\operatorname{dist}/\|(1-\Pi)X\|$, set to $0$ when the squared distance comes out
negative.}
\label{tab:dist}
\end{table}

So, to the precision available in double arithmetic,
\begin{equation}\label{eq:inExt}
 (1-\Pi)M\in\mathrm{Ext},\qquad (1-\Pi)v\in\mathrm{Ext},\qquad (1-\Pi)(M+v)\in\mathrm{Ext},
\end{equation}
with relative distance $\le1.2\times10^{-5}$ at $L=128$ and numerically zero for
$L=192,256,384$.  There is no trend with $L$: the residuals track the conditioning of
the sum \eqref{eq:closed}, not a genuine defect.  This is P1's three-way identity for
$M$ (now verified with the \emph{larger} exterior class), and it is new for $v$ and
$M+v$.

\subsection{High-precision confirmation: the distance is exactly zero}
The entries of Table~\ref{tab:dist} are differences of $O(10^{-2})$ numbers computed
through sums $\sum|\cdot|^2/d$ with $d$ down to $10^{-17}$, so ``$\pm10^{-10}$'' only says
``consistent with $0$''.  \texttt{e2\_mp.py} rebuilds the whole model in \texttt{mpmath}
from the closed forms
$P_{jk}=L^{-1}\sum_{n=0}^{L/2-1}\omega^{n(j-k)}$, $\omega=e^{2\pi i/L}$, and
$\partial_{jk}=L^{-1}\sum_n i\nu_ns_n\omega^{n(j-k)}$, and evaluates
$R:=\|M\|^2-\|\Pi M\|^2-\|\mathrm{Proj}_{\mathrm{Ext}}M\|^2$ at $60$--$80$ digits
(\texttt{e2\_mp.out}):
\begin{center}\small
\begin{tabular}{r r r r r r}
\toprule
$L$&dps&$\|M\|^2$&$\|\Pi M\|^2$&$R$&$\operatorname{dist}/\|(1-\Pi)M\|$\\\midrule
32&60&$0.055991454966480020$&$0.012011528991976492$&$-6.49\mathrm{e}{-60}$&$1.2\mathrm{e}{-29}$\\
48&60&$0.037908078979136406$&$0.007889626661341902$&$-2.46\mathrm{e}{-56}$&$9.0\mathrm{e}{-28}$\\
64&80&$0.033776977394262449$&$0.007408224258961964$&$+6.72\mathrm{e}{-72}$&$1.6\mathrm{e}{-35}$\\
\bottomrule
\end{tabular}
\end{center}
$R$ vanishes to $55$--$70$ digits and its residue tracks the working precision (it drops by
$12$ orders when \texttt{dps} goes from $60$ to $80$), not the smallest $d^c$ (which is
$10^{-16}$ at $L=48$ and $10^{-21}$ at $L=64$).  The ``leak'' of $\eta$ onto the $d=0$
modes is $<10^{-118}$ and the leak of $K$ is identically $0$.  Hence
$\operatorname{dist}((1-\Pi)M,\mathrm{Ext})=0$ \emph{exactly}, as Lemma~\ref{lem:dim}
requires, and the double-precision residuals of Table~\ref{tab:dist} are pure roundoff.

\subsection{The matrix truncation $E_ID_\chi E_I$ is not admissible}
For completeness, and as a warning, the same table for $X=P^\perp(E_ID_\chi E_I)P$:
\begin{center}\small
\begin{tabular}{rrrrr}
\toprule
$L$&$\|X\|^2$&$\|\Pi X\|^2$&$\|\Pi X\|^2/\|u\|^2$&$\operatorname{dist}((1-\Pi)X,\mathrm{Ext})^2$\\\midrule
128&$2.2819\mathrm{e}{+1}$&$1.9194\mathrm{e}{+1}$&$2608.4$&$-1.2\mathrm{e}{-13}$\\
192&$5.3598\mathrm{e}{+1}$&$4.4998\mathrm{e}{+1}$&$6188.1$&$-8.7\mathrm{e}{-13}$\\
256&$9.7335\mathrm{e}{+1}$&$8.1642\mathrm{e}{+1}$&$11051.9$&$-1.0\mathrm{e}{-12}$\\
384&$2.2368\mathrm{e}{+2}$&$1.8741\mathrm{e}{+2}$&$25290.5$&$-1.4\mathrm{e}{-11}$\\
\bottomrule
\end{tabular}
\end{center}
The invisible part is again in $\mathrm{Ext}$ (as it must be, by Lemma~\ref{lem:dim}), but
the \emph{visible} part diverges like $L^2$: sharply truncating a first-order differential
operator to a block creates an $O(L)$ boundary term with an $O(L^2)$ visible defect.  Reading
``$Z_I=sD_\chi|_I$'' as the matrix compression $E_I(sD_\chi)E_I$ therefore gives a bound
$25290\times$ worse than the compression itself, not a bound at all.  The correct reading is
the one of \S\ref{ssec:column}: the prescribed object is the column $ZE_I$, i.e.\ the full
$sD_\chi$ up to its $E_c\cdot E_c$ block, which is free anyway.

\begin{verdictgap}
Statement (L1)+(L2) and the reduction ``(UH) at leading order $\iff(1-\Pi)P^\perp Z_IP\in\mathrm{Ext}$''
are confirmed, but the Phase 4 brief's parametrisation $Z=Z_I\oplus Z_e$ must be replaced by
``$ZE_I$ prescribed, $E_cZE_c$ free''.  The conclusion is unaffected (both descriptions
give the same $\mathrm{UB}$, since $E_cD_\chi E_c$ is itself admissible), but a write-up
that literally uses $E_ID_\chi E_I$ would be wrong by four orders of magnitude.
\end{verdictgap}


\section{The optimised Uhlmann bound}\label{sec:ub}

Script \texttt{e2\_dist\_ub.py}.  By \eqref{eq:UB} and \eqref{eq:inExt},
\begin{equation}
 \mathrm{UB}_c=\tfrac12 s^2\bigl(\|u\|^2+d_c^2\bigr),\qquad
 \mathrm{UB}_{\rm rot}=\tfrac12 s^2\bigl(\|u+\eps_*v_{G_*}\|^2+d_{\rm rot}^2\bigr)
 =\tfrac12 s^2\bigl(\|u\|^2(1-\theta)+d^2_{\rm rot}\bigr),
\end{equation}
with $d_c,d_{\rm rot}$ the distances of Table~\ref{tab:dist}.  $\mathrm{UB}$ is exactly
quadratic in $s$, so all ratios are $s$-independent to machine precision; this was
confirmed at $s=0.001,0.002,0.005$ (identical to 8 digits, see \texttt{e2\_dist\_ub.out}).

\begin{table}[htbp]\centering\small
\begin{tabular}{r r r r r r}
\toprule
$L$&$\theta$&$\dfrac{\mathrm{UB}_c}{\frac12 s^2\|u\|^2}$&$\dfrac{\mathrm{UB}_{\rm rot}}{\mathrm{UB}_c}$&$1-\theta$&$\dfrac{\|\Pi(M+v)\|^2}{\|u\|^2}$\\[4pt]
\midrule
128&$0.29110$&$1.00000000$&$0.70885189$&$0.708903$&$0.708852$\\
192&$0.28880$&$0.99999999$&$0.71125039$&$0.711204$&$0.711250$\\
256&$0.28949$&$1.00000000$&$0.71058293$&$0.710508$&$0.710583$\\
384&$0.29233$&$0.99999998$&$0.70765606$&$0.707673$&$0.707656$\\
\bottomrule
\end{tabular}
\caption{The optimised Uhlmann bound.  Target values: $1$ and $1-\theta$.  For comparison,
the fixed Gaussian-taper extension of \texttt{uhlmann\_rotated.py} gives
$\mathrm{UB}_c/(\tfrac12s^2\|u\|^2)=4.565,4.389,4.324,4.305$ at $L=96,160,256,384$ and
$\mathrm{UB}_{\rm rot}/\mathrm{UB}_c=0.95$; the whole factor $4.3$ and the whole failure to
reach $1-\theta$ were artefacts of a bad exterior extension.  Also
$\mathrm{UB}_{\rm rot}/(\tfrac12 s^2\|u\|^2)$ equals the third column times the fourth and
is $0.70885,0.71125,0.71058,0.70766$.}
\label{tab:ub}
\end{table}

\subsection{Is the infimum attained at bounded cost?}
The set-theoretic identity $\mathrm{Ext}=\mathcal N^\perp$ would be worthless if the
minimiser had $\|Z_e\|_2\to\infty$.  Tikhonov regularisation ($d^c\mapsto d^c+\mu$, which
is exactly $\min\{\|X+P^\perp Z_eP\|^2+\mu\|Z_e\|_2^2\}$) gives the trade-off:

\begin{center}\small
\begin{tabular}{r rr rr rr rr}
\toprule
&\multicolumn{2}{c}{$L=128$}&\multicolumn{2}{c}{$L=192$}&\multicolumn{2}{c}{$L=256$}&\multicolumn{2}{c}{$L=384$}\\
\cmidrule(lr){2-3}\cmidrule(lr){4-5}\cmidrule(lr){6-7}\cmidrule(lr){8-9}
$\mu$&ratio&$\|Z_e\|_2/s$&ratio&$\|Z_e\|_2/s$&ratio&$\|Z_e\|_2/s$&ratio&$\|Z_e\|_2/s$\\
\midrule
$10^{-2}$ &$1.126074$&$0.541$&$1.131395$&$0.545$&$1.130633$&$0.546$&$1.132027$&$0.548$\\
$10^{-4}$ &$1.003310$&$1.220$&$1.003548$&$1.239$&$1.003570$&$1.246$&$1.003634$&$1.255$\\
$10^{-6}$ &$1.000062$&$2.069$&$1.000069$&$2.114$&$1.000070$&$2.135$&$1.000072$&$2.160$\\
$10^{-8}$ &$1.000001$&$3.031$&$1.000001$&$3.133$&$1.000001$&$3.175$&$1.000001$&$3.222$\\
$10^{-10}$&$1.000000$&$4.122$&$1.000000$&$4.283$&$1.000000$&$4.346$&$1.000000$&$4.417$\\
$10^{-12}$&$1.000000$&$5.370$&$1.000000$&$5.530$&$1.000000$&$5.632$&$1.000000$&$5.731$\\
\bottomrule
\end{tabular}
\end{center}
(``ratio'' $=\mathrm{UB}_c(\mu)/(\tfrac12 s^2\|u\|^2)$; $\|Z_e\|_2$ is the
Hilbert--Schmidt norm.)  The cost is \emph{uniformly bounded in $L$}: at $\mu=10^{-8}$ one
is within $10^{-6}$ of the infimum with $\|Z_e\|_2/s\in[3.03,3.22]$ for $L=128,\dots,384$,
growing by $6\%$ over a factor $3$ in $L$; at $\mu=10^{-4}$ one is within $0.4\%$ with
$\|Z_e\|_2/s\approx1.25$.  For scale, $\|sD_\chi\|_2/s\approx L^{1/2}\cdot O(1)$ and
$\|M\|_2/s=0.179$, so the optimal exterior generator is of the same order as the interior
one.  The infimum is therefore \emph{effectively attained}, not approached along a
divergent sequence, and there is no sign of the ill-conditioned bulk of $\mathrm{Ext}$
(\S\ref{sec:dims}) being needed.

\begin{verdictok}
$\mathrm{UB}_c=\tfrac12 s^2\|u\|^2$ and $\mathrm{UB}_{\rm rot}=(1-\theta)\,\tfrac12 s^2\|u\|^2$
to $7$--$8$ significant digits, at every $L\in\{128,192,256,384\}$ and every
$s\in\{0.001,0.002,0.005\}$, with an exterior generator of $L$-uniformly bounded norm.
At leading order (UH) holds, and with it $E_{\rm rec}\le(1-\theta)f_2z^2(1+o(1))$.
\end{verdictok}


\section{Beyond first order: exact $-\tfrac12\log\det(1-V^*V)$}\label{sec:exact}

Script \texttt{e2\_exact.py}, output \texttt{e2\_exact.out}.  With the optimised exterior
generator $Z_e$ (Tikhonov $\mu$) we form three unitaries of $\mathbb C^L$, all equal to
first order in $s$ and pairwise different at second order, and evaluate the \emph{exact}
Uhlmann exponent
\begin{equation}
 \mathrm{UH}(\widetilde W)=-\tfrac12\log\det\bigl(1-V^*V\bigr)
 =-\tfrac12\sum_j\log\bigl(1-\sigma_j^2\bigr),\qquad V=P^\perp\widetilde WP,
\end{equation}
from the singular values $\sigma_j$ of $V$ using \texttt{log1p} (Hilbert--Schmidt
quantities, double precision suffices: $\sigma_j\sim10^{-5}$ are computed with absolute
error $\sim\eps\|V\|\sim10^{-18}$, so the sum $\sim10^{-9}$ carries $\sim13$ correct
digits).  Writing $Z_{\rm pre}=-sD_\chi$ (compression) or $-s(D_\chi+G_*)$ (rotated, with
$G_*=\eps_*x$, $\eps_*=-1$):
\begin{align}
 \text{sum:}\quad&\widetilde W=\exp(Z_{\rm pre}+Z_e),&
 \text{prod:}\quad&\widetilde W=\exp(Z_{\rm pre})\exp(Z_e),&
 \text{$W_s$:}\quad&\widetilde W=W_s\,e^{-s\eps_*G_*}\exp(Z_e),
\end{align}
$W_s=\exp(-sD_\chi)$.  The first is the one-parameter-group construction for which
\texttt{rate\_of\_remainder.tex} gives the exact inequality
$\|P^\perp e^{Z}P\|_2^2\le\|P^\perp ZP\|_2^2$; the third is literally
``$W_se^{-\eps G_*}$ extended by $\exp(Z_e)$ on $I^c$''.

\begin{table}[htbp]\centering\small
\begin{tabular}{r r | r r r | r r r r}
\toprule
&&\multicolumn{3}{c|}{compression, $/\frac12 s^2\|u\|^2$}&\multicolumn{4}{c}{rotated, $/\frac12 s^2\|u\|^2$}\\
$L$&$s$&UB&sum&prod&UB&sum&prod&$W_s$\\
\midrule
128&0.001&1.00000&1.00000&0.99956&0.70885&0.70885&0.70596&0.70810\\
128&0.002&1.00000&0.99999&0.99917&0.70885&0.70884&0.70322&0.70741\\
128&0.005&1.00000&0.99993&0.99829&0.70885&0.70878&0.69599&0.70568\\
192&0.001&1.00000&1.00000&0.99950&0.71125&0.71125&0.70736&0.71053\\
192&0.002&1.00000&0.99998&0.99905&0.71125&0.71123&0.70373&0.70988\\
192&0.005&1.00000&0.99990&0.99801&0.71125&0.71110&0.69452&0.70829\\
256&0.001&1.00000&1.00000&0.99946&0.71058&0.71057&0.70576&0.70988\\
256&0.002&1.00000&0.99998&0.99897&0.71058&0.71054&0.70134&0.70925\\
256&0.005&1.00000&0.99987&0.99783&0.71058&0.71033&0.69063&0.70773\\
384&0.001&1.00000&0.99999&0.99943&0.70766&0.70764&0.70114&0.70698\\
384&0.002&1.00000&0.99997&0.99891&0.70766&0.70757&0.69542&0.70638\\
384&0.005&1.00000&0.99981&0.99770&0.70766&0.70714&0.68320&0.70499\\
\bottomrule
\end{tabular}
\caption{Exact $-\frac12\log\det(1-V^*V)$ with the optimised exterior extension,
$\mu=10^{-8}$, in units of $\frac12 s^2\|u\|^2$.  Compare $1-\theta=0.70890$, $0.71120$,
$0.71051$, $0.70767$ for $L=128,192,256,384$.  $\|Z_e\|_2/s\approx3.0$--$3.2$ and
$\|Z_e\|_{\rm op}/s\approx1.9$--$2.0$, so the exponentials are small perturbations; the
$\mu$-dependence is negligible (the $\mu=10^{-4}$ rows differ in the fourth digit, the
$\mu=10^{-6},10^{-10}$ rows in the sixth).}
\label{tab:exact}
\end{table}

Three things to read off.
\begin{enumerate}[leftmargin=1.6em,itemsep=2pt]
\item The exact value never exceeds the first-order $\mathrm{UB}$: the deficit is
 $O(s^2)$ relative ($1-0.99981=1.9\times10^{-4}$ at $s=0.005$, $1.5\times10^{-5}$ at
 $s=0.002$, a clean $s^2$ law), the expected negative quartic term of
 $-\frac12\sum\log(1-\sigma^2)=\frac12\sum\sigma^2+\frac14\sum\sigma^4+\dots$ minus the
 larger positive gain from $\|P^\perp e^ZP\|_2^2<\|P^\perp ZP\|_2^2$.
\item The rotated exponent is $(1-\theta)\cdot\frac12 s^2\|u\|^2$ to five digits at
 $s=0.001$ and slightly \emph{below} it at larger $s$.  (UH) holds with room to spare.
\item The two second-order-inequivalent constructions are both \emph{better} than the
 group version: $\exp(Z_{\rm pre})\exp(Z_e)$ gains a further $0.3$--$2.4\%$ and
 $W_se^{-s\eps_*G_*}\exp(Z_e)$ a further $0.1$--$0.4\%$.  So the choice of second-order
 completion is not a threat to (UH); it is an extra, unexploited, margin.
\end{enumerate}

\begin{verdictok}
Beyond first order (UH) holds in this model: $-\log F(\omega,\omega^{\beta_\eps})\le
2\,\mathrm{UH}(\widetilde W)$ with $\mathrm{UH}(\widetilde W)/(\tfrac12 s^2\|u\|^2)\le1-\theta$
for every $L$, $s$ and every one of the three constructions.
\end{verdictok}


\section{Verdict}\label{sec:verdict}

\subsection{Answers to the four questions of track E2}
\begin{enumerate}[leftmargin=1.8em,itemsep=3pt]
\item[(a)] \textbf{$\mathrm{Ext}=\mathcal N^\perp$ exactly, at every $L$.}  Lemma~\ref{lem:dim}:
 $\dim\mathcal N=m^2-2z^2$, $\dim\mathcal N^\perp=(L-m)^2=\dim\mathrm{Ext}$, $z=m-L/2$.
 Verified by explicit bases at $L=32,48$ (ranks exact, all principal angles $0$) and by the
 closed-form Gram spectra up to $L=384$.  The \emph{conditioning} degrades with $L$ (the
 lower quartile of the gains $\sqrt{d^c_{ij}/2}$ falls from $10^{-2.5}$ at $L=96$ to
 $10^{-7.5}$ at $L=384$), so density rather than equality is what one should expect to
 prove in the continuum.
\item[(b)] \textbf{$(1-\Pi)M$, $(1-\Pi)v$ and $(1-\Pi)(M+v)$ all lie in $\mathrm{Ext}$}, with
 relative distance $\le1.2\times10^{-5}$ at $L=128$ and $0$ (to double precision) at
 $L=192,256,384$; no trend in $L$ (Table~\ref{tab:dist}).
\item[(c)] \textbf{$\mathrm{UB}_c/(\tfrac12 s^2\|u\|^2)=1.0000000$} and
 \textbf{$\mathrm{UB}_{\rm rot}/\mathrm{UB}_c=1-\theta$} to $4$--$5$ digits, for all
 $L\in\{128,192,256,384\}$ and $s\in\{0.001,0.002,0.005\}$ ($s$-independent to machine
 precision, $\mathrm{UB}$ being exactly quadratic).  The value $4.3$ of
 \texttt{uhlmann\_rotated.py} was entirely an artefact of the fixed Gaussian-taper
 extension.  The optimum is reached with $\|Z_e\|_2/s\approx3$, uniformly in $L$.
\item[(d)] \textbf{The exact $-\frac12\log\det(1-V^*V)$ agrees} with the first-order
 $\mathrm{UB}$ to $O(s^2)$ relative and always from \emph{below}; the rotated value is
 $(1-\theta)\tfrac12 s^2\|u\|^2$ to five digits (Table~\ref{tab:exact}).
\end{enumerate}

\subsection{Consequence for Phase 4}
In this model, and to the accuracy of double-precision linear algebra,
\begin{equation}
 -\log F\bigl(\omega,\omega^{\beta_{\eps_*}}\bigr)\ \le\ (1-\theta)\,f_2z^2\,(1+o(1)),
 \qquad \theta\approx0.29,
\end{equation}
so the upper half of the Phase 4 target holds and $E_{\rm rec}\le(1-\theta)f_2z^2(1+o(1))$:
\emph{the zero-collar geometric compression is not optimal in exact fidelity}.  Moreover,
since $\mathrm{Ext}=\mathcal N^\perp$ (not merely $\mathrm{Ext}\ni(1-\Pi)M$), the same
conclusion holds for \emph{every} isometric channel $\mathrm{Ad}\,\Gamma(W_I)$, and the
second-order recovery problem over the isometric orbit is exactly the $g_Q$ problem.

\subsection{What is not settled}
\begin{enumerate}[leftmargin=1.8em,itemsep=2pt]
\item These are finite-$L$ statements.  The continuum question is whether
 $\overline{\mathrm{Ext}}=\mathcal N^\perp$ for the Hardy projection on an arc, and
 Lemma~\ref{lem:dim} does not survive as a dimension count.  The evidence that it does
 survive is quantitative, not structural: the required $\|Z_e\|_2/s$ is bounded
 ($3.03,3.13,3.18,3.22$ at $L=128,192,256,384$; a $6\%$ rise over a factor $3$ in $L$).
 A proof is track E1's; the natural target is not (Q) itself but
 $(1-\Pi)M\in\overline{\mathrm{Ext}}$ with a norm bound, which is what the numerics
 actually support.
\item[$\checkmark$] (Closed.)  The double-precision distances of (b) were re-run in
 \texttt{mpmath} at $60$--$80$ digits for $L=32,48,64$: $\operatorname{dist}((1-\Pi)M,\mathrm{Ext})^2$
 vanishes to $55$--$70$ digits, with a residue that tracks the working precision.  The
 same was not done at $L\ge128$ (cost), but there is no mechanism by which the identity
 could fail there given Lemma~\ref{lem:dim}.
\item The brief's parametrisation $Z=Z_I\oplus Z_e$ is not literally correct for this
 $\chi$ (\S\ref{ssec:column}); the value of $\mathrm{UB}$ is unaffected, but a write-up
 using the matrix truncation $E_ID_\chi E_I$ would be wrong by a factor $2.5\times10^4$
 at $L=384$.
\item Everything is at one geometry ($a=2$, $L_g=1$) and one UV taper ($0.30,0.42$);
 \texttt{s10\_theta.py}'s \texttt{diagnostics()} shows $\theta$ is stable under both, but
 the $\mathrm{UB}$ ratios were not re-scanned over them.
\end{enumerate}

\begin{verdictok}
Track E2 is complete and affirmative on all four items.  The single quantitative headline:
with an optimised exterior extension the Uhlmann bound for the compression \emph{equals}
$\frac12 s^2\|u\|^2=g_Q(\delta_c)$ exactly (not $4.3\times$ it), and for the S10 rotated
channel it equals $(1-\theta)$ times that, with $\theta=0.29$ --- confirmed at the level of
the exact Fredholm determinant, not only at second order.
\end{verdictok}


\end{document}
